\cotefile{156-6}
\scanpage{42}\mspage{\deleted{\unsure{39|34}} 41}

\deleted{le sous-atelier $\unsure{\mathcal F|\mathcal F'}$ à partir de $\mathcal F$, comme} \deleted{\unsure{l'ens des} $F$}
\[
\deleted{\mathcal F'=\{F\in\mathcal F\mid\widetilde F\subset M'\}\mathrel{\overset{\mathrm{dfn}}{=}}\mathrm{Env}(M')}
\]
\deleted{Donc} \deleted{On trouve ainsi}

\uline{Proposition} (Ne dépend pas en fait \unsure{d'une} donnée $\ll$ sur $\mathcal F$, \inserted{seulemt de $\leq$}) On a deux applications \unsure{biunivoques} inverses l'une de l'autre
\diagram{\begin{tikzcd}[ampersand replacement=\&,column sep=large,row sep=small]
{\substack{\deleted{\unsure{\mathrm{Satsp}|\mathrm{Satpr}}}\\\mathrm{Sousatspac}(\mathcal F)}} \arrow[r,shift left=1ex,"{\mathcal F'\longmapsto\mathcal F'\cap M}"] \& \mathrm{Ferm}(M) \arrow[l,shift left=1ex] \\
{\{F\in\mathcal F\mid\widetilde F\subset M\}\mathrel{\overset{\mathrm{dfn}}{=}}\mathrm{Env}(M')} \& M' \arrow[l,maps to]
\end{tikzcd}}
% slanted
\margin{\deleted{\unsure{Dans le cas} $\mathcal F$} Si on identifie $\mathcal F$ à une partie de $\unsure{\mathcal P|\mathrm P}(M)$, alors $\mathrm{Env}.(M')=\{F\in\mathcal F\mid F\subset M'$}
entre l'ensemble des \deleted{\unsure{sous-ateliers}} \inserted{parties \unsure{propres|fermées}} spacieuses \inserted{$\mathcal F'$} de $\mathcal F$, et \unsure{l'ens} des parties \uline{fermées} \inserted{$M'$} de $M$.

Il faut voir, \deleted{que si} \inserted{quand $\mathcal F$ est un atelier, si} $\mathcal F'$ est un sous-atelier de $\mathcal F$, \deleted{\unsure{alors}} pour $\leq$, $\ll$, \deleted{\unsure{c'est un}} \inserted{au sens des} ateliers $\mathrm{At}\,1$ -- $\mathrm{At}\,7$. Rappelons pour mémoire $\mathrm{At}\,1$ -- $\mathrm{At}\,3$

\uline{At 1} \deleted{\unsure{Les parties} \illegible{1}} \unsure{Les} sup majorés \inserted{(pour $\leq$)} \unsure{existent} dans $\mathcal F$ ; et $\mathcal F\ne\varnothing$ (ou encore, a un plus petit élément $\varnothing_{\mathcal F}$).

\uline{At} 2 Toute \deleted{\unsure{fig}} $F\in\mathcal F$ est borne sup. pour $\leq$ des ``figures élémentaires'' (\unsure{les|des} \unsure{pts|élts|elts} \unsure{strictmt.} irr.)\unsure{.} qu'elle majore (pour $\leq$)

\uline{At} 3 Si $F,G\in\mathcal F$, pour que $F\vee\unsure{C|G}$ existe dans $\mathcal F$, \inserted{(\unsure{i.e.|i.e.,} $\{F,G\}$ majoré)} il suffit que \deleted{\unsure{si}} pour $X\mathrel{\triangleleft}F$, $Y\mathrel{\triangleleft}C$, $X\vee Y$ existe.
\scanpage{43}\mspage{42}

Pour At 1, At 2 c'est clair, pour At 3 il faut \deleted{de plus} imposer, en plus de

\uline{Sousat 1} $\mathcal F'$ est une partie fermée \inserted{non vide} de $\mathcal F$,

la condition

\uline{Sousat} 2 Si $F,G\in\mathcal F'$ et si pour \deleted{\unsure{tout}} $X\in\widetilde F$, $Y\in\widetilde G$, $X\vee Y$ existe dans $\mathcal F'$ (i.e. \deleted{$X\mathrel{\substack{\scriptscriptstyle\times\\[-.7ex]\scriptscriptstyle\times}}Y$\textsuperscript{\unsure{$*$|$\times$}} \textsuperscript{$*$}dans $\mathcal F$, alors $X\vee Y$ existe \unsure{donc|dans}, $X\vee Y\in\mathcal F'$ et $X\mathrel{\substack{\scriptscriptstyle\times\\[-.7ex]\scriptscriptstyle\times}}Y\unsure{\notin|\in}\mathcal F'$, alors $X\mathrel{\substack{\scriptscriptstyle\times\\[-.7ex]\scriptscriptstyle\times}}Y$ dans $\mathcal F$ et $X\vee Y\in\mathcal F'$} $X,Y$ majorés \unsure{ds|dans} $\mathcal F'$) alors $F\vee\unsure{C|G}\in\mathcal F'$ (NB on sait que $F\vee C$ existe dans $\mathcal F$).
% slanted
\margin{\uline{NB} C'est une condition \unsure{intrinsèque} à $(\mathcal F',\leq)$, i.e. la condition At 3 pour $\mathcal F'$.}

Les conditions At 4, At 5, At 6, At 7 \inserted{(pages 21, 27)} passent trivialement \inserted{une à une} de $\mathcal F$ à $\mathcal F'$. Pour At 7, je note seulement en passant
\[
\begin{aligned}
&\text{Si }X,Y\in M',\ \text{alors }X\mathrel{\overset{\circ}{\ll}}Y\text{ dans }\mathcal F\text{ équivaut à}\\
&\hspace{12em}X\mathrel{\overset{\circ}{\ll}}Y\text{ dans }\mathcal F'.
\end{aligned}\tag*{\makebox[0pt][r]{\makebox[\displaywidth][l]{(\unsure{$*$|M})}}}
\]

\deleted{Supposons maintenant $\mathcal F$ interprété comme un ensemble de figures (en fait, un sous-atelier) de $\mathrm{Fig}(M)$.}

\uline{Définition} Soit $\mathcal F$ un atelier au sens At 1 -- At 7. On appelle \uline{sous-atelier} de $\mathcal F$, une partie \deleted{\unsure{fermée}}\inserted{$\mathcal F'$} de $\mathcal F$ \deleted{\unsure{sat}}, satisfaisant \deleted{\unsure{à}}\unsure{.} Sousat 1 et Sousat 2 ci-dessus, i.e. telle que $1^\circ)$ $\mathcal F'$ soit une partie fermée non vide de $(\mathcal F,\leq)$ \deleted{et} (ce qui n'est autre que Sousat 1) et $2^\circ)$ $\mathcal F'$, munie \unsure{de de|de} $\leq$, $\ll$, est un atelier (ce qui équivaut à Sous-at 2), On dit que
\scanpage{44}\mspage{43}

$\mathcal F'$ est un \uline{sous-atelier} \uline{spacieux}, \unsure{si l'est|si c'est} une partie fermée \uline{spacieuse} de $\mathcal F$ \deleted{(\unsure{auquel cas} la condition Sousat 2 est} Ce qui suffit à entrainer que c'est un sous-atelier.

\uline{Exemple} Soit $\mathcal F$ un atelier, et considérons l'application canonique \unsure{:|.}
\[
\mathrm{Mulomb}\colon\mathcal F\hookrightarrow\mathrm{Fig}(M),
\]
on a vu que l'image est une partie fermée pour $\leq$, et que les relations $\leq$ et $\ll$ sur $\mathcal F$ sont induites par celles de $\mathrm{Fig}(M)$. Ainsi $\mathcal F$ est un sous-atelier de l'atelier des figures \unsure{dans} $M$. \deleted{En fait, on a vu aussi que c'est} \inserted{Mais ce n'est pas en général} un sous-atelier spacieux (mais ça l'est peut-être dans le cas ensembliste, à voir).

\par\noindent\hfill\rule{2cm}{.3pt}\hfill\mbox{}\par

\deleted{Je reviens à la \unsure{m} aux ``petites'' \unsure{ombres} de}

\uline{Remarque} Soit $\mathcal F$ un ordonné pour $\leq$, $\mathcal F'$ une \unsure{partie} fermée \inserted{non vide} de $\mathcal F$. Alors les parties fermées non vides de $\mathcal F'$ sont \unsure{exactmt|exactt} les parties fermées non vides de $\mathcal F$ qui sont contenues dans $\mathcal F'$. On a le même énoncé pour les parties fermées\unsure{.} spacieuses, pourvu que l'on suppose $\mathcal F'$ spacieux. De façon précise, si on a des
\scanpage{45}\mspage{44}

inclusions
\[
\mathcal F''\subset\mathcal F'\subset\mathcal F^{\deleted{\unsure{\prime|\prime\prime}}}
\]
de parties fermées non vides, si $\mathcal F''$ est spacieux dans $\mathcal F'$, il l'est dans $\mathcal F'$, et la réciproque est vraie si $\mathcal F'$ est spacieux dans $\mathcal F$. De $\overline{\mathrm m}$, si $\mathcal F$ est un atelier, $\mathcal F'$ un sous-atelier, les sous-ateliers de $\mathcal F'$ sont les sous-ateliers de $\mathcal F$ contenus dans $\mathcal F''$ -- en d'autres termes, la condition supplémentaire Sousat 2 pour $\mathcal F''$ est la $\overline{\mathrm m}$, qu'on la formule pour $\mathcal F''\subset\mathcal F'$ ou pour $\mathcal F''\subset\mathcal F$ ; en effet, \unsure{c'est} une condition intrinsèque à $\mathcal F''$ (savoir, At 3 pour $\mathcal F''$). Donc, \deleted{\unsure{le}}
\[
\mathrm{Sousat}(\mathcal F')=\{\mathcal F''\in\mathrm{Sousat}(\mathcal F)\mid\mathcal F''\subset\mathcal F'\}
\tag*{\makebox[0pt][r]{\makebox[\displaywidth][l]{(1,50)}}}
\]
et si $\mathcal F'$ spacieux dans $\mathcal F$
\[
\mathrm{Sousatspac}(\mathcal F')=\{\mathcal F''\in\mathrm{Sousatspac}(\mathcal F)\mid\mathcal F''\subset\mathcal F'\}.
\tag*{\makebox[0pt][r]{\makebox[\displaywidth][l]{(1\unsure{;|,}51)}}}
\]

\uline{Remarque} Soit $(\mathcal F,\leq)$ un un ordonné satisfaisant At 1, -- At \unsure{2|3}, \deleted{\unsure{donc}} qui se décrit donc par $(M,\leq,\mathcal F)$, \deleted{$\mathcal F$ \unsure{étant}} $(M\leq)$ ensemble ordonné, $\mathcal F\deleted{\unsure{\in}}\subset\unsure{\mathcal P|\mathrm P}\unsure{f|{}_{f}}(M)$ \inserted{ens. des parties \uline{fermées} de $M$} telle que $1^\circ)$ $\lvert\mathcal F\rvert=M$ i.e. $\forall x\in M$, $\exists F\in\mathcal F$ avec $X\in F$, et $2^\circ)$ $\mathcal F$ fermé dans $\mathcal P f(M)$ i.e. si $F\in\mathcal F$ et $G$ partie fermée de $F$, alors $G\in\mathcal F$. \unsure{Prenons|Posons} dans $M$ la relation $\ll\mathrel{\overset{\mathrm{dfn}}{=}}\leq$. Alors, si $(M,\leq)$ satisfait At 3, alors $(M,\leq,\ll)$ satisfait At 1 à At 7. \unsure{Parmi|Pour mi}
\scanpage{46}\mspage{45}

axiomes qui vont suivre, il y en a qui vont éliminer, je pense, cet exemple trivial d'atelier.

\par\noindent\hfill\rule{1.4cm}{.3pt}\hfill\mbox{}\par

Revenons aux \uline{petites} ombres et multiombres. C'est le moment d'introduire enfin

\uline{At 8} Soit $X\in M$, et soit $\unsure{\delta|\partial} X=\mathop{\mathrm{Sup}_{\deleted{\unsure{\mathrm{au}|\mathrm{aux}}}}}\limits_{\substack{Y\mathrel{\unsure{\triangleleft|\trianglelefteq}}X\\Y\ne X}}^{(\leq)}\,Y$.

Alors $\exists\,x\in\unsure{\mathcal L|L}$ tel que $x\unsure{\ll|\leq}X$, \deleted{$x\unsure{\in|<}$} $x\mathrel{\not\ll}\delta X$ \deleted{\unsure{et}} (i.e. $x\mathrel{\overset{\circ}{\ll}}X$). \deleted{et} \deleted{\unsure{ou un} En}

En d'autres termes, posant \deleted{\unsure{p}}
\[
\begin{aligned}
\mathrm{omb}(X)^\circ&=\mathrm{omb}(X)-\mathrm{omb}(\delta X)\\
&=\mathrm{omb}(X)-\bigcup_{\substack{Y\mathrel{\triangleleft} X\\Y\ne X}}\mathrm{omb}(Y)\unsure{,}\ ,
\end{aligned}
\]
on a :
\[
\boxed{\mathrm{omb}(X)^\circ\ne\varnothing}
\tag*{\makebox[0pt][r]{\makebox[\displaywidth][l]{(1.52)}}}
\]

Considérons alors, pour $F\in\mathcal F$, \deleted{\unsure{sa}} multiombre, qui est par definition \deleted{formée des \unsure{part} parties} \inserted{la figure \unsure{ens.iste|ensembliste}}
\[
\mathrm{mulomb}(F)=\{\mathrm{omb}(X)\mid\underset{\inserted{\text{i.e. }X\in\widetilde F}}{X\mathrel{\triangleleft} F},\ \mathrm{omb}(X)^\circ\ne\varnothing\}.
\]
At 8 nous dit que l'ens d'indices à prendre ici est $\widetilde F$ tout entier. Ainsi on trouve \inserted{(cf \unsure{sorite|sorte} ensembliste)}

\uline{Corollaire} \deleted{\unsure{Question|Soient}} $F$ une figure, $X,Y\in\widetilde F$. Alors conditions suivantes équivalentes
\[
\unsure{\text{a)}|\text{.)}}\quad X\leq Y\qquad(\Longleftrightarrow X\ll Y)
\]
\scanpage{47}\mspage{46}

\[
\begin{aligned}
\text{b)}\quad&\mathrm{omb}(X)\subset\mathrm{omb}(Y)\\
\text{c)}\quad&\mathrm{omb}(X)^\circ\subset\mathrm{omb}(Y)\\
\text{d)}\quad&\mathrm{omb}(X)^\circ\cap\mathrm{omb}(Y)\ne\varnothing
\end{aligned}
\]

Il y a une correspondance \unsure{biunivoque}
\diagram{\begin{minipage}{\linewidth}
\[
\vcenter{\hbox{\begin{tikzcd}[ampersand replacement=\&,column sep=large,row sep=small]
{} \arrow[d,phantom,"{\left\{\vphantom{\dfrac{\displaystyle F}{\displaystyle F}}\right.}" description] \& F' \arrow[r,maps to] \& \mathrm{mulomb}(F')\\
{} \& \mathrm{Ssfig}(F) \arrow[r,"\sim"] \arrow[d,phantom,"{\scriptscriptstyle\Vert}" description] \& \mathrm{Ssfig}(\mathrm{mulomb}(F))\\
\& \mathcal F_{\leq F} \&
\end{tikzcd}}}
\tag*{\makebox[0pt][r]{\makebox[\displaywidth][l]{(1.53)}}}
\]
\end{minipage}}

\uline{Corollaire \unsure{1|I}} Pour que $F$ soit élémentaire \unsure{i.e.|resp.} une \unsure{multistrate}, il f. et s. que $\mathrm{mulomb}(F)$ le soit.

\uline{Remarque} Il n'y a aucune raison que At 8 soit stable par passage à un sous-atelier $\mathcal F'$, $\overline{\mathrm m}$ quand celui-ci est spacieux. Mais \deleted{considérons} \inserted{supposons que} $\deleted{\unsure{G|F}}F\in\mathcal F'$, $G\in\mathcal F$, $\unsure{C|G}\ll F$ implique $C\ll F$ i.e. que $\mathcal F'$ soit fermé pour $\ll$ également -- on dit alors que le sous-atelier est \uline{plein}. En général, on aura
\[
\unsure{\mathcal L|L}\cap\mathcal F'\subset\mathcal L'\quad(\mathrel{\overset{\mathrm{dfn}}{=}}\text{ ens des lieux de }\mathcal F')
\]
mais dans le cas d'un atelier plein$^{*}$, on aura
\[
\mathcal L\cap\mathcal F'=\mathcal L'.
\]
% slanted
\margin{il suffit qu'il soit $\mathcal L$-plein, i.e. que $F\in\mathcal F'$, $x\in\mathcal L$, $x\ll F\Rightarrow x\in\mathcal F'$}
Dans ce cas$^{*}$, si $\mathcal F$ satisfait At 8, \deleted{$\unsure{L|\mathcal L}$}$\mathcal F'$ aussi. De plus, \deleted{\unsure{si}} pour $F\in\mathcal F'$, $\mathrm{omb}(F)$ et $\mathrm{mulomb}(F)$ sont les $\overline{\mathrm m}$s, au sens de $\mathcal F$ ou au sens de $\mathcal F'$.
\scanpage{48}\mspage{47}

\deleted{\uline{Corollaire 2} Soient $F,\unsure{C|G}\in\mathcal F$, avec $G<F$. Alors}
\[
\deleted{\exists\,\unsure{,}\,x\in\mathrm{omb}(F),\ \text{t.q.}}
\]
\deleted{\unsure{J'arrive|J'en viens|Je reviens|J'arriverai} à la notion d'atelier \uline{ensembliste}} \deleted{\unsure{Déf}}

Je reprends les \unsure{sorites} des \deleted{homo}morphismes d'\unsure{Ateliers|ateliers}

\uline{Définition} Soient $\mathcal F',\deleted{\unsure{G|\mathcal G}}\mathcal F$ deux ateliers. On appelle \deleted{\unsure{prémo}} \uline{quasi-morphisme} \deleted{de $\mathcal F$} (en attendant une notion plus forte de morphisme !) de $\mathcal F'$ dans $\mathcal F$, une application
% slanted
\margin{il vaudrait peut-être mieux de parler de \uline{co}morphismes (\unsure{évt.|evt.} ``quasi''), cf plus bas.}
\[
f:\mathcal F'\longrightarrow\mathcal F
\]
qui est

a) croissante pour $\leq$ et $\ll$

b) commute aux Sup \deleted{\unsure{même} (\unsure{majorés|compris})} (et transforme \inserted{donc} $\varnothing_{\mathcal F'}$ en $\varnothing_{\mathcal F}$) et aux Inf pour $\leq$

On n'exige pas qu'elle transforme fig. élémentaire en figure élémentaire. Mais quand il en est ainsi, cette application est connue quand on connait sa restriction
\[
f_{\unsure{\mathrm{str}1|\mathrm{str}l}}:M'\unsure{\mathrel{\relbar\joinrel\relbar\joinrel\shortmid}|\longrightarrow}M
\]
aux ens. de \unsure{multistrates} \inserted{(= fig. élém.)}. Dans ce cas, pour que $f$ soit \unsure{injective|injection}, il f. et s. que $f_{\mathrm{str}1}^{\unsure{F|\mathcal F}}$ le soit (et en termes des plongements $\mathcal F\hookrightarrow\unsure{\mathcal P|\mathrm P}(M)$, $\mathcal F'\unsure{\lhook\joinrel\relbar|\hookrightarrow}\mathcal P(M')$, $f$ s'identifie alors à la restriction \deleted{\unsure{/corestriction}} de \deleted{$\unsure{f}$} $A\to f_{\mathrm{str}1}(A)$ : $\mathcal P(M')\longrightarrow\mathcal P(M)$, du moins quand $f(M')$ est une partie \inserted{$\leq$-fermée de $M$, cf plus bas.}
% slanted
\margin{pas clair, il \deleted{\unsure{pour}} faudrait que \inserted{$F'\in\mathcal F'$}, on ait $\widetilde{f(F')}=f(\widetilde{F'})$, ce qui n'est pas clair, cf plus bas.}

Il faudrait arriver à imposer aux quasi-morph. des conditions \unsure{supplt|suppl}, qui impliquent que lorsque $f$ est injectif, alors l'application induite :
\scanpage{49}\mspage{48}

\diagram{\begin{minipage}{\linewidth}
\[
\vcenter{\hbox{\begin{tikzcd}[ampersand replacement=\&,column sep=large,row sep=small]
\mathrm{Ssfig}(F') \arrow[r] \& \mathrm{Ssfig}(F^{\unsure{9|\prime|?}}) \& \text{où }F=f(F') \\
G' \arrow[r,maps to] \& f(\unsure{C|G}') \&
\end{tikzcd}}}
\tag*{\makebox[0pt][r]{\makebox[\displaywidth][l]{(1.54)}}}
\]
\end{minipage}}
% slanted
\margin{\unsure{ce qui|qui} implique que \deleted{\unsure{Il est sté}} \deleted{\unsure{que}} $f(\mathcal F')$ \deleted{\unsure{est}}soit une partie fermée de $\mathcal F$ pour $\leq$}
soit \uline{bijective}. Sans supposer $f$ injective, dans tous les cas à ma connaissance, on aura une application canonique (\unsure{croissante})
\diagram{\begin{minipage}{\linewidth}
\[
\vcenter{\hbox{\begin{tikzcd}[ampersand replacement=\&,row sep=large]
{\deleted{\unsure{\psi|q}}}\ \& \widetilde F \arrow[d,"{\unsure{\varphi|\phi}_{f,F'}=\varphi}"] \\
\& \widetilde F'
\end{tikzcd}}}
\tag*{\makebox[0pt][r]{\makebox[\displaywidth][l]{(1.55)}}}
\]
\end{minipage}}
d'ens ordonnés, de telle façon que (1.54) en soit déduite par l'application en sens inverse
\diagram{\begin{minipage}{\linewidth}
\[
\vcenter{\hbox{\begin{tikzcd}[ampersand replacement=\&,column sep=large,row sep=small]
{} \arrow[d,phantom,"{\left\{\vphantom{\dfrac{\displaystyle F}{\displaystyle F}}\right.}" description] \& {\unsure{\mathcal P|\mathrm P}\unsure{f|{}_{f}}(\widetilde F')} \arrow[r,dash,"{\unsure{\shortmid|>}}"{at end,description,inner sep=0pt}] \& \mathcal P f(\widetilde F)\\
{} \& A \arrow[r,maps to] \& \varphi^{-1}(A)\ ,
\end{tikzcd}}}
\tag*{\makebox[0pt][r]{\makebox[\displaywidth][l]{(1.56)}}}
\]
\end{minipage}}
et de plus, la \deleted{topologie de} relation d'ordre sur $\widetilde F$ est image inv. de celle sur $\widetilde F'$.

Dans le cas où $f$ est injective, et où la condition (1.54) est satisfaite, (1.55) sera un \unsure{izo.|iso.} d'ens. ordonnés. Dans le cas général, on \unsure{déduira|désirera} encore que l'application (1\unsure{:|.}54) est \deleted{\unsure{tout}} au moins \uline{surjective}. \deleted{\unsure{Le}} Mais dans des cas impor\unsure{\deleted{tants}|tants} (comme $\mathcal F\to\mathrm{Fig}(\unsure{\mathcal L|L})$), même si $f$ n'est pas injective, on aura (1.54) bijective pour \unsure{tte|toute} $F$.

On notera que dans les cas où (1.54) peut se déduire d'un \deleted{\unsure{ap}} hom d'ens ordonnés (1.55), celui-ci est unique. En d'autres termes
